% GENERATED FILE. Edit canonical lessons, then run npm run generate:books.
% canonical-source-hash: fnv1a64:ea9c9b8139641db6
% canonical-lessons: FR-C385-bureaucratie, FR-C385-demarche, FR-C385-procedure, FR-C385-certificat, FR-C385-attestation

\chapter{La paperasse}
\label{ch:fr-la-paperasse}
\hlchaptermodality{385}

\begin{chapteropening}
\textbf{By the end of this chapter:} \emph{I can deal with red tape: bureaucracy, the steps to take, a procedure, a certificate and an official statement.}

Complete the last lesson of chapter 385: I can deal with red tape: bureaucracy, the steps to take, a procedure, a certificate and an official statement.
\end{chapteropening}

\section[la bureaucratie]{la bureaucratie — the bureaucracy, the red tape}

\label{lesson:FR-C385-bureaucratie}



\begin{quote}
Before the new one: say the French for to renew, then the French for the renewal.
\end{quote}

\subsection*{You'll want to know: la bureaucratie}
\textbf{la bureaucratie} — \textquotedblleft{}the bureaucracy, the red tape\textquotedblright{}.

\begin{quote}
Il y a trop de bureaucratie dans ce pays.
\end{quote}

\begin{quote}
\emph{There is too much red tape in this country.}
\end{quote}

The \emph{t} is said like \emph{s}: \emph{bureaucrassie}.

\begin{cousinweb}
\emph{bureaucratie} is a French formation from \emph{bureau} \textquotedblleft{}office\textquotedblright{} and \emph{-cratie} \textquotedblleft{}rule\textquotedblright{} (EtymDB), from Greek \emph{krátos} \textquotedblleft{}power\textquotedblright{}: rule by offices.
\end{cousinweb}

\begin{grammarlens}[title={What you've built}]
One more word for telling, explaining and describing at B1.
\end{grammarlens}

\subsection*{Your turn}
\begin{itemize}
\raggedright
  \item \emph{Say it:} \emph{la bureaucratie}
  \item \emph{Say it:} \emph{la bureaucratie}, once more
  \item \emph{Say it:} a sentence of your own with \emph{la bureaucratie}
  \item \emph{Recall:} say the French for to renew, then the French for the renewal, then say \emph{la bureaucratie} again
\end{itemize}

\begin{tcolorbox}[breakable,colback=teal!4,colframe=teal!35!black,title={Before you move on}]
What does la bureaucratie mean? (\textquotedblleft{}The bureaucracy, the red tape\textquotedblright{}.) Say it once more.
\end{tcolorbox}
\section[la démarche]{la démarche — the step (to get something done); the gait}

\label{lesson:FR-C385-demarche}



\begin{quote}
Before the new one: say the French for the renewal, then the French for the bureaucracy.
\end{quote}

\subsection*{You'll want to know: la démarche}
\textbf{la démarche} — \textquotedblleft{}the step (to get something done); the gait\textquotedblright{}.

\begin{quote}
J'ai fait toutes les démarches pour obtenir mon visa.
\end{quote}

\begin{quote}
\emph{I took all the steps to get my visa.}
\end{quote}

\emph{Faire des démarches} — to take the necessary steps. Also the way someone walks.

\begin{cousinweb}
\emph{démarche} is formed on an old verb \emph{démarcher} (EtymDB), built from \emph{dé-} and \emph{marcher} \textquotedblleft{}to walk\textquotedblright{}; \emph{marcher} is usually traced to a Frankish word for treading.
\end{cousinweb}

\begin{grammarlens}[title={What you've built}]
One more word for telling, explaining and describing at B1.
\end{grammarlens}

\subsection*{Your turn}
\begin{itemize}
\raggedright
  \item \emph{Say it:} \emph{la démarche}
  \item \emph{Say it:} \emph{la démarche}, once more
  \item \emph{Say it:} a sentence of your own with \emph{la démarche}
  \item \emph{Recall:} say the French for the renewal, then the French for the bureaucracy, then say \emph{la démarche} again
\end{itemize}

\begin{tcolorbox}[breakable,colback=teal!4,colframe=teal!35!black,title={Before you move on}]
What does la démarche mean? (\textquotedblleft{}The step (to get something done); the gait\textquotedblright{}.) Say it once more.
\end{tcolorbox}
\section[la procédure]{la procédure — the procedure}

\label{lesson:FR-C385-procedure}



\begin{quote}
Before the new one: say the French for the bureaucracy, then the French for the step.
\end{quote}

\subsection*{You'll want to know: la procédure}
\textbf{la procédure} — \textquotedblleft{}the procedure\textquotedblright{}.

\begin{quote}
La procédure est longue, mais simple.
\end{quote}

\begin{quote}
\emph{The procedure is long, but simple.}
\end{quote}

In law, a court case: \emph{une procédure judiciaire}.

\begin{cousinweb}
\emph{procédure} is built on \emph{procéder} \textquotedblleft{}to proceed\textquotedblright{}, from Latin \emph{prōcēdere} \textquotedblleft{}to go forward\textquotedblright{} (Lewis \& Short), built on \emph{cēdere} \textquotedblleft{}to go\textquotedblright{}.
\end{cousinweb}

\begin{grammarlens}[title={What you've built}]
One more word for telling, explaining and describing at B1.
\end{grammarlens}

\subsection*{Your turn}
\begin{itemize}
\raggedright
  \item \emph{Say it:} \emph{la procédure}
  \item \emph{Say it:} \emph{la procédure}, once more
  \item \emph{Say it:} a sentence of your own with \emph{la procédure}
  \item \emph{Recall:} say the French for the bureaucracy, then the French for the step, then say \emph{la procédure} again
\end{itemize}

\begin{tcolorbox}[breakable,colback=teal!4,colframe=teal!35!black,title={Before you move on}]
What does la procédure mean? (\textquotedblleft{}The procedure\textquotedblright{}.) Say it once more.
\end{tcolorbox}
\section[le certificat]{le certificat — the certificate}

\label{lesson:FR-C385-certificat}



\begin{quote}
Before the new one: say the French for the step, then the French for the procedure.
\end{quote}

\subsection*{You'll want to know: le certificat}
\textbf{le certificat} — \textquotedblleft{}the certificate\textquotedblright{}.

\begin{quote}
Le médecin m'a donné un certificat pour l'école.
\end{quote}

\begin{quote}
\emph{The doctor gave me a certificate for school.}
\end{quote}

\emph{Un certificat médical} — a doctor's note.

\begin{cousinweb}
\emph{certificat} comes from Medieval Latin \emph{certificātum} (EtymDB), built on Latin \emph{certus} \textquotedblleft{}sure\textquotedblright{} and \emph{facere} \textquotedblleft{}to make\textquotedblright{} (Lewis \& Short): a paper that makes a fact sure.
\end{cousinweb}

\begin{grammarlens}[title={What you've built}]
One more word for telling, explaining and describing at B1.
\end{grammarlens}

\subsection*{Your turn}
\begin{itemize}
\raggedright
  \item \emph{Say it:} \emph{le certificat}
  \item \emph{Say it:} \emph{le certificat}, once more
  \item \emph{Say it:} a sentence of your own with \emph{le certificat}
  \item \emph{Recall:} say the French for the step, then the French for the procedure, then say \emph{le certificat} again
\end{itemize}

\begin{tcolorbox}[breakable,colback=teal!4,colframe=teal!35!black,title={Before you move on}]
What does le certificat mean? (\textquotedblleft{}The certificate\textquotedblright{}.) Say it once more.
\end{tcolorbox}
\section[l'attestation]{l'attestation — the official statement, the attestation}

\label{lesson:FR-C385-attestation}



\begin{quote}
Before the new one: say the French for the procedure, then the French for the certificate.
\end{quote}

\subsection*{You'll want to know: l'attestation}
\textbf{l'attestation} — \textquotedblleft{}the official statement, the attestation\textquotedblright{}.

\begin{quote}
La banque demande une attestation de domicile.
\end{quote}

\begin{quote}
\emph{The bank is asking for proof of address.}
\end{quote}

A signed paper that confirms a fact. Feminine.

\begin{cousinweb}
\emph{attestation} is a learned borrowing from Late Latin \emph{attestātiō} (EtymDB), from \emph{attestārī} \textquotedblleft{}to bear witness to\textquotedblright{} (Lewis \& Short), built on \emph{testis} \textquotedblleft{}witness\textquotedblright{}.
\end{cousinweb}

\begin{grammarlens}[title={What you've built}]
One more word for telling, explaining and describing at B1.
\end{grammarlens}

\subsection*{Your turn}
\begin{itemize}
\raggedright
  \item \emph{Say it:} \emph{l'attestation}
  \item \emph{Say it:} \emph{l'attestation}, once more
  \item \emph{Say it:} a sentence of your own with \emph{l'attestation}
  \item \emph{Recall:} say the French for the procedure, then the French for the certificate, then say \emph{l'attestation} again
\end{itemize}

\begin{tcolorbox}[breakable,colback=teal!4,colframe=teal!35!black,title={Before you move on}]
What does l'attestation mean? (\textquotedblleft{}The official statement, the attestation\textquotedblright{}.) Say it once more.
\end{tcolorbox}
